% Generated draft: requires mathematical and source review.
\documentclass[11pt]{article}
\usepackage[margin=25mm]{geometry}
\usepackage{fontspec}
\setmainfont{Latin Modern Roman}
\usepackage{amsmath,amssymb,mathtools,unicode-math}
\setmathfont{Latin Modern Math}
\usepackage{enumitem,needspace,xurl,hyperref}
\setlength{\parindent}{0pt}
\setlength{\parskip}{5pt}
\setlength{\emergencystretch}{3em}
\newcommand{\N}{\mathbb{N}}
\newcommand{\Z}{\mathbb{Z}}
\newcommand{\Q}{\mathbb{Q}}
\newcommand{\R}{\mathbb{R}}
\newcommand{\C}{\mathbb{C}}
\newcommand{\F}{\mathbb{F}}
\newcommand{\A}{\mathbb{A}}
\newcommand{\PP}{\mathbb{P}}
\newcommand{\E}{\mathbb{E}}
\newcommand{\eps}{\varepsilon}
\newcommand{\dd}{\,\mathrm{d}}
\begin{document}
\section*{UnsolvedMath ID: 701328}
\textbf{Source problem code:} catalogue entry
\subsection{Definitions and mathematical statement}
A right order on a group \(G\) is a strict total order \(<_G\) satisfying
\[
 g<_G h\ \Longrightarrow\ gt<_G ht
 \qquad(g,h,t\in G).
\]
A group is right-orderable if it admits a right order. A right-ordered group is a pair \((G,<_G)\) with a specified right order. An order-preserving embedding \(\iota:(G,<_G)\hookrightarrow(H,<_H)\) is an injective group homomorphism satisfying \(g<_G h\Longrightarrow\iota(g)<_H\iota(h)\).

An equation over \(G\) has the displayed form
\[
 w(a_1,\ldots,a_k,x_1,\ldots,x_n)=1,
\]
where the \(a_i\) are specified elements of \(G\), the \(x_j\) are unknown group elements, and \(w\) is a group word. A solution in an overgroup consists of a group embedding \(\iota:G\hookrightarrow H\) and elements \(b_1,\ldots,b_n\in H\) such that
\[
 w\bigl(\iota(a_1),\ldots,\iota(a_k),b_1,\ldots,b_n\bigr)=1_H.
\]

\textbf{Unresolved scope of equations.} The supplied record does not specify which equations are eligible for the proposed closure property. In particular, it supplies no hypothesis excluding equations already incompatible with injectivity on the coefficient group, and it does not say whether solvability in some group overgroup is assumed. An unrestricted quantifier over all group words would include such incompatible equations. Thus a precise unconditional formulation cannot be recovered from the supplied text alone.

To display the requested closure conditions without silently adding a hypothesis, write \(\mathcal E(G)\) for the unspecified family of equations over \(G\) intended by the source. This is a placeholder for missing source information. The two questions are then:
\begin{enumerate}
\item For every right-orderable group \(G\) and every equation \(w=1\) in \(\mathcal E(G)\), do there exist a right-orderable group \(H\), a group embedding \(\iota:G\hookrightarrow H\), and a solution of \(w=1\) in \(H\) under \(\iota\)?
\item For every right-ordered group \((G,<_G)\) and every equation \(w=1\) in \(\mathcal E(G)\), do there exist a right-ordered group \((H,<_H)\), an order-preserving group embedding \(\iota:(G,<_G)\hookrightarrow(H,<_H)\), and a solution of \(w=1\) in \(H\) under \(\iota\)?
\end{enumerate}
The second question quantifies over the specified right order on \(G\), rather than merely over the existence of some right order.

The source reports an affirmative result for equations with a single constant, \(k=1\), citing \emph{J. Group Theory} \textbf{11} (2008), 623--633. This supplied assertion is retained with its source qualification; its precise hypotheses and its coverage of the two closure questions require checking the cited article.
\subsection{Short English statement}
The problem asks whether equations over a right-orderable group can be solved in a right-orderable overgroup, and whether a specified right order can also be preserved by the embedding. The supplied record reports an affirmative single-constant case, but does not specify the admissible equations, so the general formulation remains conditional on clarifying that scope.
\subsection{Sources}
\begin{itemize}
\item Local UnsolvedMath registry: permanent ID \(701328\), exact source problem code \texttt{catalogue entry}, title \emph{Kourovka Notebook Problem 17.26}; supplied upstream revision \texttt{d371ee917ac9647fbb06cd6fc529deec6f3dda82}.
\item \emph{The Kourovka Notebook (Unsolved Problems in Group Theory)}, Problem 17.26, proposed by V. V. Bludov and A. M. W. Glass, 17th issue (2010). The supplied record also identifies the 21st issue edition (2026), \texttt{arXiv:1401.0300v48}, and identifies E. I. Khukhro and V. D. Mazurov as editors of the collection.
\item Supplied reference for the affirmative single-constant case: \emph{J. Group Theory} \textbf{11} (2008), 623--633. The record supplies no article title, authors, DOI, or link for this reference.
\item The registry labels the problem \texttt{open} and records an update timestamp of \texttt{2026-10-16T14:00:00Z}. These are supplied registry data, not an independently verified mathematical or literature assessment. No source was independently consulted for this draft.
\end{itemize}
\end{document}
